\documentclass[11pt]{article}
\usepackage{geometry}
\geometry{margin=1in}
\usepackage{graphicx}
\graphicspath{ {./final_figures/} }
\usepackage{float}

\title{Baseline Vitamin D, Preeclampsia,\\
and Early-Pregnancy Gene Expression}
\author{Chunlei Li}
\date{}

\begin{document}
\maketitle

\section{Introduction}
Low maternal vitamin D in early pregnancy has been linked to a higher
risk of preeclampsia \cite{Mirzakhani2016}. CYP24A1 encodes the enzyme
that inactivates the active form of vitamin D, so I examined its
expression together with baseline vitamin D and pregnancy outcome. I also
examined UTS2 and RUNDC3A, the two most variable of the supplied genes,
to compare expression patterns. This report is descriptive and does not
test hypotheses.

\section{Methods}

\subsection{Data source and sample size}
The data come from the study by Mirzakhani et al.\ \cite{Mirzakhani2016}
and are available from the Gene Expression Omnibus under accession
GSE85307 \cite{GEO_GSE85307}. The course provided three files: covariate
data for 157 participants, expression values for 500 genes measured with
the Affymetrix Human Gene 1.0 ST Array, and a linker file mapping probe
IDs to gene names. Blood samples were collected between 10 and 19 weeks of
gestation. The initial sample size was 157 participants (110 normal
pregnancies and 47 with preeclampsia).

\subsection{Data processing and final sample size}
All data wrangling was performed in R. I removed the text label from each
covariate value, converted the continuous covariates to numbers, matched
expression columns to participants by their GEO sample ID, and replaced
probe IDs with gene names using the linker file. Every sample ID and probe
ID matched, and no covariate or expression value was missing. I therefore
removed no participants, and the final sample size was 157. I derived a
categorical vitamin D group from baseline 25-hydroxyvitamin D using a
cutoff of 30~ng/mL ($<$30 vs.\ $\geq$30~ng/mL).

\subsection{Analyses}
I inspected histograms of each continuous covariate before choosing
summary statistics. Maternal age was approximately symmetric and is
reported as mean (SD); vitamin D and BMI were right-skewed and are
reported as median [Q1, Q3]. Categorical variables are reported as
n (\%). The summary table is stratified by pregnancy outcome.

For CYP24A1, UTS2, and RUNDC3A, I used one plotting function to produce a
histogram of expression, a scatterplot of expression against baseline
vitamin D with an unadjusted linear fit and 95\% confidence band, and a
boxplot of expression by pregnancy outcome and vitamin D group. For the
heatmap, I selected the ten genes with the highest standard deviation
across participants and applied hierarchical clustering with Euclidean
distance and complete linkage to rows and columns, without row
standardization. As an additional plot, I used a dumbbell plot to compare the median
expression of the same ten genes between the two outcome groups, with
genes ordered by the difference in medians (preeclampsia minus normal). All
analyses were descriptive; I did not perform hypothesis tests or estimate
adjusted effects.

\subsection{Software}
Analyses were conducted in R version 4.6.1 \cite{RCoreTeam}.
Figures were made with ggplot2 version 4.0.3
\cite{Wickham2016} and pheatmap version 1.0.13
\cite{Kolde_pheatmap}. Tables were made with knitr version
1.52 \cite{Xie2015} and kableExtra version
1.4.1 \cite{Zhu_kableExtra}, and the analysis was
compiled with R Markdown \cite{Xie2018}.

\subsection{Code availability}
All code is available in my public GitHub repository:\\
https://github.com/llliii0088979-a11y/QBS103\_DataScience\_Project.

\subsection{Use of AI and other resources}
This was my first time using R, so I made frequent errors while writing
the code. I used an AI assistant (Claude) to help me debug errors,
suggest code for parts of the analysis (such as the dumbbell plot), and
draft and edit the wording of this report. I also watched YouTube
tutorials to help me understand R and ggplot2 concepts. I chose the
research question, genes, and covariates, reviewed and ran all of the
code myself, and checked that the results and text matched the data.

\section{Results}

\subsection{Participant characteristics}
I found lower median baseline vitamin D (17.7 vs.\ 21.8~ng/mL) and
higher median BMI (31.24 vs.\ 26.54~kg/m$^2$) in participants with
preeclampsia than in participants with normal pregnancies, while mean
maternal age was similar (25.78 vs.\ 26.63 years). A larger share of the
preeclampsia group had vitamin D below 30~ng/mL (87.2\% vs.\ 70.9\%;
Table~\ref{tab:one}).

\begin{table}[H]
\centering
\caption{Participant characteristics by pregnancy outcome.}
\label{tab:one}
\footnotesize
\begin{tabular}{|l|r|r|r|}
\hline
 & Overall & Normal & Preeclampsia \\
Variable & ($n = 157$) & ($n = 110$) & ($n = 47$) \\
\hline
Vitamin D (ng/mL), median [Q1, Q3] & 20.4 [13.4, 29.5] & 21.8 [13.85, 32.08] & 17.7 [12.55, 27.4] \\
Maternal age (years), mean (SD) & 26.38 (5.04) & 26.63 (5.08) & 25.78 (4.95) \\
Maternal BMI (kg/m$^2$), median [Q1, Q3] & 27.33 [22.9, 33.27] & 26.54 [22.35, 30.97] & 31.24 [24.62, 36.41] \\
\hline
Pregnancy outcome, n (\%) & & & \\
Normal & 110 (70.1) & 110 (100) & 0 (0) \\
Preeclampsia & 47 (29.9) & 0 (0) & 47 (100) \\
\hline
Vitamin D group, n (\%) & & & \\
$<$30 ng/mL & 119 (75.8) & 78 (70.9) & 41 (87.2) \\
$\geq$30 ng/mL & 38 (24.2) & 32 (29.1) & 6 (12.8) \\
\hline
\end{tabular}
\end{table}

\subsection{CYP24A1}
I observed a single, narrow peak in CYP24A1 expression near 4.8, with a
small upper tail (Figure~\ref{fig:cyp-hist}). The fitted line between
baseline vitamin D and CYP24A1 expression was nearly flat, so I did not
see a clear linear pattern in this sample (Figure~\ref{fig:cyp-scatter}).
Expression distributions overlapped across all four outcome and vitamin D
subgroups. I interpreted the lower median in the preeclampsia,
$\geq$30~ng/mL subgroup cautiously because it included only six
participants (Figure~\ref{fig:cyp-box}).

\begin{figure}[H]
  \centering
  \includegraphics[width=0.65\textwidth]{cyp-hist-1.pdf}
  \caption{CYP24A1 expression in all 157 participants. The histogram uses
  24 bins and the supplied expression scale.}
  \label{fig:cyp-hist}
\end{figure}

\begin{figure}[H]
  \centering
  \includegraphics[width=0.65\textwidth]{cyp-scatter-1.pdf}
  \caption{CYP24A1 expression and baseline vitamin D. The line is an
  unadjusted linear fit; shading shows the 95\% confidence interval for
  the fitted mean.}
  \label{fig:cyp-scatter}
\end{figure}

\begin{figure}[H]
  \centering
  \includegraphics[width=0.65\textwidth]{cyp-box-1.pdf}
  \caption{CYP24A1 expression by pregnancy outcome and vitamin D group.
  Boxes show medians and quartiles; whiskers extend to observations within
  1.5 times the IQR, and points beyond them are outliers. Group sizes:
  $<$30~ng/mL, Normal 78 and Preeclampsia 41; $\geq$30~ng/mL, Normal 32
  and Preeclampsia 6.}
  \label{fig:cyp-box}
\end{figure}

\subsection{UTS2}
UTS2 expression was clearly bimodal: most participants had values near
5, and a smaller group had values between about 6 and 8.5
(Figure~\ref{fig:uts2-hist}). Both groups appeared across the range of
baseline vitamin D, so the single linear fit passed between them and did
not describe either group well (Figure~\ref{fig:uts2-scatter}). The six
preeclampsia participants with vitamin D $\geq$30~ng/mL had a higher
median UTS2 expression, but I did not treat this small subgroup as
evidence of a difference (Figure~\ref{fig:uts2-box}).

\begin{figure}[H]
  \centering
  \includegraphics[width=0.65\textwidth]{uts2-hist-1.pdf}
  \caption{UTS2 expression in all 157 participants. The histogram uses 24
  bins and the supplied expression scale.}
  \label{fig:uts2-hist}
\end{figure}

\begin{figure}[H]
  \centering
  \includegraphics[width=0.65\textwidth]{uts2-scatter-1.pdf}
  \caption{UTS2 expression and baseline vitamin D. The line is an
  unadjusted linear fit; shading shows the 95\% confidence interval for
  the fitted mean.}
  \label{fig:uts2-scatter}
\end{figure}

\begin{figure}[H]
  \centering
  \includegraphics[width=0.65\textwidth]{uts2-box-1.pdf}
  \caption{UTS2 expression by pregnancy outcome and vitamin D group.
  Boxes show medians and quartiles; whiskers extend to observations within
  1.5 times the IQR, and points beyond them are outliers. Group sizes:
  $<$30~ng/mL, Normal 78 and Preeclampsia 41; $\geq$30~ng/mL, Normal 32
  and Preeclampsia 6.}
  \label{fig:uts2-box}
\end{figure}

\subsection{RUNDC3A}
I observed a broader, roughly symmetric distribution of RUNDC3A
expression centered near 9 (Figure~\ref{fig:rundc3a-hist}). The fitted
line sloped slightly downward with baseline vitamin D, with substantial
scatter around it (Figure~\ref{fig:rundc3a-scatter}). Median RUNDC3A
expression was lower in the preeclampsia group within both vitamin D
groups, although the distributions overlapped
(Figure~\ref{fig:rundc3a-box}).

\begin{figure}[H]
  \centering
  \includegraphics[width=0.65\textwidth]{rundc3a-hist-1.pdf}
  \caption{RUNDC3A expression in all 157 participants. The histogram uses
  24 bins and the supplied expression scale.}
  \label{fig:rundc3a-hist}
\end{figure}

\begin{figure}[H]
  \centering
  \includegraphics[width=0.65\textwidth]{rundc3a-scatter-1.pdf}
  \caption{RUNDC3A expression and baseline vitamin D. The line is an
  unadjusted linear fit; shading shows the 95\% confidence interval for
  the fitted mean.}
  \label{fig:rundc3a-scatter}
\end{figure}

\begin{figure}[H]
  \centering
  \includegraphics[width=0.65\textwidth]{rundc3a-box-1.pdf}
  \caption{RUNDC3A expression by pregnancy outcome and vitamin D group.
  Boxes show medians and quartiles; whiskers extend to observations within
  1.5 times the IQR, and points beyond them are outliers. Group sizes:
  $<$30~ng/mL, Normal 78 and Preeclampsia 41; $\geq$30~ng/mL, Normal 32
  and Preeclampsia 6.}
  \label{fig:rundc3a-box}
\end{figure}

\subsection{Ten-gene heatmap}
Because I kept the supplied scale, the colors mainly reflected
differences in overall expression level between genes, and the genes
clustered into a lower-expression group (UTS2, CEACAM6, CLIC2, ZNF594) and
a higher-expression group. One cluster of participants had higher UTS2
expression, matching the second UTS2 group in Figure~\ref{fig:uts2-hist}.
The tracking bars showed no visible separation by pregnancy outcome or
vitamin D group (Figure~\ref{fig:heatmap}). I treated the clustering as
exploratory.

\begin{figure}[H]
  \centering
  \includegraphics[width=\textwidth]{heatmap-1.pdf}
  \caption{The ten genes with the highest standard deviation across 157
  participants. Colors show expression on the supplied scale. Rows and
  columns use Euclidean distance and complete-linkage hierarchical
  clustering. Tracking bars show pregnancy outcome and baseline vitamin D
  group. Sample labels are hidden for readability.}
  \label{fig:heatmap}
\end{figure}

\subsection{Median expression of the ten heatmap genes by outcome}
Among the ten most variable genes, I found the largest difference
between outcome groups for RUNDC3A, with lower median expression in preeclampsia
(8.74 vs.\ 9.15; difference $-0.41$), followed by TGM2 ($-0.28$), ZNF594
($-0.21$), and CLIC2 ($-0.19$). The remaining six genes differed by 0.06
or less (Figure~\ref{fig:dumbbell}). I report these as unadjusted
descriptive differences, not test results.

\begin{figure}[H]
  \centering
  \includegraphics[width=0.8\textwidth]{dumbbell-1.pdf}
  \caption{Median expression of the ten highest-variance genes in normal
  pregnancies ($n = 110$) and preeclampsia ($n = 47$). Each line joins the
  two group medians for one gene; numbers on the right give the difference
  (preeclampsia minus normal) on the supplied expression scale. Genes are
  ordered by this difference.}
  \label{fig:dumbbell}
\end{figure}

\section{Discussion}
In this sample, participants with preeclampsia had lower baseline vitamin
D and were more often below 30~ng/mL, consistent with the direction reported by
Mirzakhani et al.\ \cite{Mirzakhani2016}. CYP24A1 expression showed no
clear pattern with vitamin D or outcome in my analysis. Among the most variable genes,
RUNDC3A had the largest outcome difference, consistent with its lower
median in the preeclampsia group within both vitamin D groups. I consider these findings descriptive:
the analyses are unadjusted, include no hypothesis tests, and some
subgroups are small, especially the six preeclampsia participants with
vitamin D $\geq$30~ng/mL.

\bibliographystyle{unsrt}
\bibliography{references}

\end{document}
